% Uniform phase signs for the Eulerian-square polynomial.
\section{Uniform phase signs in the central region}\label{sec:bulk}

Throughout this section $n\ge250$. We use the real-rootedness of the
factorial-weighted polynomial
\[
 F_n(t)=\sum_{j=1}^n j!A(n,j)t^j
       =n!t\prod_{i=1}^{n-1}(t+\beta_i),\qquad \beta_i>0.
\]
For real $s$, set $C_s=1+e^{-s}$ and
\[
 \mathcal R_{n,s}(z)=\frac{F_n(C_s/z)}z,\qquad
 z_q=s+iq\pi\quad(q\text{ odd}).
\]
The pole expansion and its continuous nearest-pole phase are
\[
 P_n(-e^{-s})=C_s\sum_{q\text{ odd}}\mathcal R_{n,s}(z_q),
 \qquad
 \phi(s)=\arg\mathcal R_{n,s}(s+i\pi).
\]
The sum is absolutely convergent. All phase comparisons below use the
continuous argument specified by the factorization of $F_n$.

\begin{lemma}[Central phase signs]\label{lem:bulk}
For every $n\ge250$ and $|s|\le4\sqrt n$,
\[
 \sum_{q=3,5,7,\ldots}
 \frac{|\mathcal R_{n,s}(s+iq\pi)|}
      {|\mathcal R_{n,s}(s+i\pi)|}<1.
\]
In particular, if $\phi(s)=j\pi$ in this interval, then
\[
 \operatorname{sgn}P_n(-e^{-s})=(-1)^j.
\]
One may use the upper bounds $.85$ for $0\le s\le3\sqrt n$,
$.33$ for $3\sqrt n\le s\le4\sqrt n$, and $.88$ for
$-4\sqrt n\le s\le0$.
\end{lemma}

We prove the three estimates separately. The finite arithmetic checks
used in the proof involve only rational constants; their precise
inequalities and bounding series are given at the end of the section.

\subsection{A weighted mean and the positive half of the interval}

\begin{lemma}\label{lem:weightedmean}
For $n\ge2$ and $t\ge e/n$,
\[
 \frac{tF_n'(t)}{F_n(t)}\ge\frac{n+1}{2}.
\]
\end{lemma}
\begin{proof}
Pair the terms indexed by $j$ and $\ell$ with $j+\ell=n+1$ and
$j>\ell$. Eulerian symmetry gives $A(n,j)=A(n,\ell)$. The geometric
mean of $\ell+1,\ldots,j$ is at least $(n!)^{1/(n-1)}$: these intervals
have the same center, and deleting an outer symmetric pair increases
the average logarithm, since the products of the symmetric pairs
increase toward the center. Moreover,
\[
 \log(n!)\ge\int_1^n\log x\,dx
 =n\log n-n+1\ge(n-1)(\log n-1).
\]
Thus $j!t^j\ge\ell!t^\ell$. Each pair makes a nonnegative contribution
to $\sum_j(j-(n+1)/2)j!A(n,j)t^j$, proving the assertion.
\end{proof}

Fix $s\ge0$, put $C=C_s$ and $a_i=C/\beta_i$, and write
\[
 R_q=\frac{|\mathcal R_{n,s}(s+iq\pi)|}
 {|\mathcal R_{n,s}(s+i\pi)|}.
\]
Direct factorization gives
\begin{equation}\label{eq:bulk-positive-ratio}
 R_q=\left(\frac{s^2+\pi^2}{s^2+q^2\pi^2}\right)^{(n+1)/2}
 \prod_{i=1}^{n-1}
 \left(\frac{(s+a_i)^2+q^2\pi^2}{(s+a_i)^2+\pi^2}\right)^{1/2}.
\end{equation}
For $0<s\le3\sqrt n$, $C/s\ge e/n$, so Lemma~\ref{lem:weightedmean}
implies
\[
 \sum_i\frac{a_i}{s+a_i}\ge\frac{n-1}{2}.
\]
At least $(n-1)/3$ of the $a_i$ are at least $s/3$: if $M$ have this
property, the left side is at most $M+(n-1-M)/4$.
Set
\[
 v=\frac{s^2}{\pi^2},\quad c=\frac{16}{9},\quad
 M_0=\frac{n-1}{3},\quad V=\frac{9n}{\pi^2}.
\]
Bounding the $M$ distant factors separately in
\eqref{eq:bulk-positive-ratio} gives
\begin{equation}\label{eq:positive-fq}
 R_q\le f_q(v):=\frac{v+1}{v+q^2}
 \left\{\frac{(v+1)(cv+q^2)}{(v+q^2)(cv+1)}\right\}^{M_0/2}.
\end{equation}
The sign of the logarithmic derivative of $f_q$ is the sign of
\[
 \left[c^2+\frac{M_0}{2}c(c-1)\right]v^2
 +c(1+q^2)v+q^2\left[1-\frac{M_0}{2}(c-1)\right].
\]
Its first two coefficients are positive. Hence $f_q$ decreases and
then increases, or increases throughout, and its maximum on $[0,V]$
is attained at an endpoint. At zero, $f_q(0)=q^{-2}$. Using
$\log(1-x)\le-x$ at the other endpoint gives
\[
 f_q(V)\le e^{-E_n(q)},\qquad
 E_n(q)=\frac{7n(n-1)(q^2-1)\pi^2}
 {6(9n+q^2\pi^2)(16n+\pi^2)}.
\]
For fixed $q$, $E_n(q)$ increases with $n$. At $n=250$,
\[
 E_n(3)>.60,\quad E_n(5)>1.70,\quad
 E_n(7)>3.11,\quad E_n(9)>4.60.
\]
Thus the first three ratios are less than $.55,.19,.045$, respectively,
and $f_9(V)<1/81$. Differentiation in $y=q^2$ shows that
$yf_{\sqrt y}(V)$ decreases whenever
\[
 y\ge\frac{cV}{M_0(c-1)/2-1}<13.
\]
Consequently $f_q(V)\le q^{-2}$ for $q\ge9$, and the endpoint argument
extends this inequality to every $v\in[0,V]$. Since
\[
 \sum_{q=9,11,\ldots}q^{-2}
 =\frac{\pi^2}{8}-1-\frac19-\frac1{25}-\frac1{49}<.063,
\]
we obtain $\sum_{q\ge3,\ q\text{ odd}}R_q<.848<.85$. At $s=0$ the
same result follows directly from $R_q\le q^{-2}$.

For the remaining positive interval, $3\sqrt n\le s\le4\sqrt n$,
Lemma~\ref{lem:weightedmean} still applies because
$C/s\ge1/(4\sqrt n)\ge e/n$. Put
\[
 p_i=\frac{a_i}{s+a_i},\quad t_i=1-p_i,\quad
 A=\frac{q^2\pi^2}{s^2},\quad B=\frac{\pi^2}{s^2},
\]
so $\sum_i p_i\ge(n-1)/2$. Define
\[
 h_q(t)=\frac12\log\frac{1+At^2}{1+Bt^2},\qquad D_q=h_q(1).
\]
If $A\le3$ and $B\le.005$, then $h_q(t)\le tD_q$ on $[0,1]$.
Indeed, for
$\psi(v)=v/(1+v)-\tfrac12\log(1+v)$,
\[
 t^2\left(\frac{h_q(t)}t\right)'=\psi(At^2)-\psi(Bt^2).
\]
The function $\psi$ increases on $[0,1]$ and decreases on $[1,3]$.
If $At^2\le1$, the difference is nonnegative. Otherwise
$\psi(At^2)\ge\psi(3)>.05$, whereas $\psi(Bt^2)\le B/2\le.0025$.
It follows from \eqref{eq:bulk-positive-ratio} that
\begin{equation}\label{eq:positive-strong}
 \log R_q\le-\frac{n+3}{2}D_q
 \qquad\left(\frac{q\pi}{s}\le\sqrt3\right).
\end{equation}
Let $Q$ be the largest odd integer at most $\sqrt3s/\pi$. Then
$Q\ge25$ and
\[
 \frac{Q\pi}{s}>\sqrt3\frac{Q}{Q+2}
 \ge\sqrt3\frac{25}{27}>1.6.
\]
The inequality $\log(1+x)\ge2x/(2+x)$ in
\eqref{eq:positive-strong} gives, for $q=3,5,7,9$,
\[
 -\log R_q\ge
 \frac{(q^2-1)9.86}{64+2(q^2+1)9.87/250}.
\]
The corresponding four upper bounds sum to less than $.326$.
For $11\le q\le Q$ the same inequality gives
\[
 -\log R_q\ge\alpha(q^2-1),\qquad \alpha=\frac{9.86}{160.16},
\]
and hence
\[
 \sum_{11\le q\le Q,\ q\text{ odd}}R_q
 \le\frac{e^{-120\alpha}}{1-e^{-48\alpha}}<.001.
\]
For later use, comparison of the individual root factors yields, for
any real $s$, any odd $Q$, and $q\ge Q$,
\begin{equation}\label{eq:universal-pole-tail}
 \frac{R_q}{R_Q}\le\left(\frac Qq\right)^2
 \exp\!\left(\frac{(n+1)s^2}{2Q^2\pi^2}\right),\qquad
 \sum_{q>Q,\ q\text{ odd}}R_q
 \le\frac Q2R_Q\exp\!\left(\frac{(n+1)s^2}{2Q^2\pi^2}\right).
\end{equation}
Here the last inequality uses
$\sum_{q>Q,\ q\text{ odd}}q^{-2}\le1/(2Q)$.
At the present $Q$,
$D_Q\ge\tfrac12\log(3.56/1.005)>.62$, so $R_Q<e^{-.31n}$.
The extra exponential in \eqref{eq:universal-pole-tail} is at most
$e^{.197n}$ and $Q/2<2\sqrt n$. The resulting tail is less than $.001$
for $n\ge250$. This proves the positive-interval bound $.33$.

\subsection{A uniform variance estimate}

We next record the distributional estimate needed on the negative
half of the interval. Write $(n)_\ell=n!/(n-\ell)!$.

\begin{lemma}\label{lem:variance}
Let $n\ge250$, $\pi\le u\le4\sqrt n$, and
\[
 B=\frac{1+e^u}{u},\qquad Y=\log(nB).
\]
If $L$ has probabilities
\[
 p_\ell\ \propto\ \frac{A(n,\ell+1)}{(n)_\ell B^\ell},
 \qquad 0\le\ell\le n-1,
\]
then $L$ is a sum of independent Bernoulli variables and
\[
 \operatorname{Var}L<\frac{8n}{Y^2}.
\]
\end{lemma}
\begin{proof}
The real-rootedness of $F_n$, coefficient reversal, and positive
tilting show that $L$ is a Bernoulli sum. More explicitly, its success
probabilities are $\beta_i/(B+\beta_i)$. In particular its probabilities
are log-concave. Let $m$ be a mode and put $N=n-m$.

The sequence $u_k=A(n,k)^{1/n}$, with zero endpoints, is concave.
Indeed, the $n$-dimensional volume of the section
$\sum_{i=1}^{n+1}x_i=k$ of $[0,1]^{n+1}$ is
$\sqrt{n+1}A(n,k)/n!$; Brunn's concavity theorem applies
\cite{Gardner2002}. Thus $u_{k+1}/u_k\le(k+1)/k$.
Eulerian unimodality places $m$ before the midpoint. If $m\ge1$,
\[
 B(n-m+1)\le\frac{A(n,m+1)}{A(n,m)}
 \le(1+1/m)^n\le e^{n/m}.
\]
Consequently
\begin{equation}\label{eq:bulk-mode-mean}
 m\le\frac n{Y-\log2},\qquad
 \mu:=\mathbb EL\le\frac n{Y-\log2}+1,
\end{equation}
where the second bound is Darroch's mode theorem
\cite{Darroch1964}. The first bound also holds when $m=0$.

Put $h=Y^2/n$. Elementary endpoint estimates give
\begin{equation}\label{eq:variance-parameters}
 Y>7.5,\quad Y/n<.26,\quad h<17,\quad
 N/n>.852,\quad\log(N/n)>-.16.
\end{equation}
For the two upper bounds, $Y$ increases with $u$ and at $u=4\sqrt n$,
\[
 Y=4\sqrt n+\tfrac12\log n-\log4+
       \log(1+e^{-4\sqrt n}).
\]
The resulting bounds on $Y/n$ and $Y/\sqrt n$ decrease after $n=250$;
the endpoint gives $Y<65$ and $Y/\sqrt n<\sqrt{17}$.
The lower bound follows at $n=250,u=\pi$.

Set
\[
 a=\frac{p_{m-1}}{p_m},\quad b=\frac{p_{m+1}}{p_m},\quad
 c=(BN)^{1/n}-1,\quad e=1-[B(N+1)]^{-1/n},\quad
 J=\lfloor n/Y\rfloor\ge3,
\]
with missing adjacent probabilities set equal to zero. Concavity of
$u_k$ gives the tangent estimates
\begin{align}
 \frac{p_{m+j}}{p_m}&\le b\,e^{-hj(j-1)/7}
       &&(1\le j\le J),\label{eq:variance-forward}\\
 \frac{p_{m-j}}{p_m}&\le a\,e^{-hj(j-1)/3}
       &&(1\le j\le m),\label{eq:variance-backward}\\
 ab&\le e^{-h/2}.\label{eq:variance-adjacent}
\end{align}
Here are the details retaining the adjacent factors. If
$c_0=u_{m+2}/u_{m+1}-1\le c$, concavity implies
$u_{m+j+1}/u_{m+1}\le1+jc_0$. After subtracting $\log b$, the
$c_0$-dependent expression
$n[\log(1+jc_0)-\log(1+c_0)]$ is increasing in $c_0$.
Therefore
\[
 \log\frac{p_{m+j}}{bp_m}
 \le n[\log(1+jc)-j\log(1+c)]
      +\sum_{i=0}^{j-1}-\log(1-i/N)
 \le-\left\{\frac{nc^2}{(1+Jc)(1+c)}-\frac1{N-J}\right\}
       \frac{j(j-1)}2.
\]
The elementary bounds
\[
 c\ge\left(1-\frac{.16}{7.5}\right)\frac Yn,
 \quad c<.298,\quad Jc<1.15,\quad (N-J)/n>.719
\]
make the bracket exceed $2h/7$.
For the backward direction the actual decrement
$d=1-u_m/u_{m+1}$ is at least $e$. Subtracting $\log a$ leaves
$n[\log(1-jd)-\log(1-d)]$, which decreases in $d$. Hence
\[
 \log\frac{p_{m-j}}{ap_m}
 \le n[\log(1-je)-j\log(1-e)]
       +\sum_{i=0}^{j-1}\log(1+i/(N+1))
 \le-\left(ne^2-\frac1{N+1}\right)\frac{j(j-1)}2.
\]
Since $1-e^{-t}\ge.88t$ for $0\le t\le.26$, the bracket exceeds
$2h/3$. This argument is needed only for $m\ge1$, so $N+1\le n$.
Finally, concavity gives
$u_mu_{m+2}/u_{m+1}^2\le(1-d)(1+d)\le1-e^2$, and thus
\[
 ab\le(1-e^2)^n\frac{N+1}{N}
 \le\exp\!\left(-ne^2+\log\frac{N+1}{N}\right)<e^{-h/2}.
\]
When $m=0$, the backward assertion is empty and $ab=0$.

The one-step version of the forward estimate gives
$p_{m+J}/p_{m+J-1}\le e^{-2h(J-1)/7}$.
All later ratios are no larger, by log-concavity. For $t=n/Y\ge3$,
\[
 \frac{J(J-1)}{t^2}\ge\frac38,\qquad
 \frac{J-1}{t}\ge\frac12.
\]
The initial amplitude is at most $e^{-3n/56}$ and the subsequent
ratio is less than $e^{-Y/7}<.343$. Summing the geometric tail gives
\begin{equation}\label{eq:variance-far-tail}
 \sum_{j>J}j^2\frac{p_{m+j}}{p_m}
 \le1.3J^2e^{-3n/56}
 \le1.3(n/7.5)^2e^{-3n/56}<.003.
\end{equation}

Let $v=\operatorname{Var}L$. Fourier inversion for a Bernoulli sum gives
\[
 \max_k p_k\le e^{-v}I_0(v),\qquad
 I_0(v)=\sum_{k\ge0}\frac{(v^2/4)^k}{(k!)^2}.
\]
Indeed, the squared modulus of its characteristic function is at most
$\exp[-4v\sin^2(t/2)]$. We shall use
\begin{equation}\label{eq:variance-mass}
 \sqrt v\,e^{-v}I_0(v)<.47\qquad(v\ge.7).
\end{equation}
For $v\ge1$, put $f(v)=\sqrt v\,e^{-v}I_0(v)$. The Bessel equation gives
$f''+2f'+f/(4v^2)=0$. Writing $I_1=I_0'$, the defining series imply
$2I_1(1)<I_0(1)$, so $f'(1)<0$; at every zero of $f'$, one has
$f''<0$. Thus $f$ decreases on $[1,\infty)$, and $f(1)<.47$.
On $[.7,1]$, inequality~\eqref{eq:variance-mass} is certified on
300 rational intervals of length $.001$ by the bounds below.

Define
\[
 S(t)=\sum_{j\ge1}j^2e^{-tj(j-1)},\qquad
 T(t)=\sum_{j\ge2}(j-1)^2e^{-tj(j-1)}.
\]
Equations~\eqref{eq:variance-forward}--\eqref{eq:variance-far-tail}
imply
\[
 v\le\mathbb E(L-m)^2
 \le p_m\{bS(h/7)+aS(h/3)+.003\}
 \le p_m\{S(h/7)+e^{-h/2}S(h/3)+.003\}.
\]
The last inequality maximizes a linear function over
$0\le a,b\le1$, $ab\le e^{-h/2}$, using $S(h/7)\ge S(h/3)$.
If $h\le11$ and $v\ge8/h$, then $v>.7$ and
\[
 v^{3/2}<.47\{S(h/7)+e^{-h/2}S(h/3)+.003\}
          <(8/h)^{3/2},
\]
a contradiction. For $0<h\le2$, the last comparison follows from
$S(t)\le2t^{-3/2}$ for $0<t\le2/3$. To verify this elementary bound,
apply the integral-plus-maximum bound for a positive unimodal sum to
$x^2e^{-tx(x-1)}$. After multiplication by $t^{3/2}$, its integral
increases with $t$ and is less than $1.21$ at $t=2/3$; its maximum
also increases with $t$ and is less than $.75$ there. For the remaining
values of $h\le11$, the six rational interval checks in
Table~\ref{tab:variance-intervals} prove the comparison.

For $11\le h<17$, clamp $L$ to $[m,m+1]$ if $b\ge a$, and to
$[m-1,m]$ otherwise. Write the resulting variable as $Z$, and
$W=L-Z$. In the first case $a\le e^{-h/4}$, and in the second
$b\le e^{-h/4}$. Since $\operatorname{Var}Z\le1/4$,
\[
 v\le\frac14+\mathbb EW^2+\sqrt{\mathbb EW^2}
 \le\frac14+\Delta(h)+\sqrt{\Delta(h)},
\]
where
\[
 \Delta(h)=\max\{T(h/7)+e^{-h/4}S(h/3),\,
                    T(h/3)+e^{-h/4}S(h/7)\}+.003.
\]
This function decreases. On each half-unit interval $[a,b]$ in
$[11,17]$, exact rational bounds give
\[
 8/b-1/4-\Delta(a)>0,\qquad
 \Delta(a)<[8/b-1/4-\Delta(a)]^2.
\]
Consequently $v<8/h$, completing the proof.
\end{proof}

\subsection{Kernels for the negative half of the interval}

Write $s=-u$, $u\ge0$. For $u\ge\pi$ set
\[
 B=\frac{1+e^u}{u},\quad Y=\log(nB),\quad
 r_i=\beta_i/B,\quad p_i=\frac{r_i}{1+r_i},\quad
 w_i=p_i(1-p_i).
\]
Then $\mu=\sum p_i$ and $v=\sum w_i$ are the mean and variance in
Lemma~\ref{lem:variance}. For $q\ge1$ let
\[
 D_q=\frac12\log\frac{u^2+q^2\pi^2}{u^2+\pi^2},\qquad
 g_q(r)=\frac12\log
 \frac{u^2(1-1/r)^2+q^2\pi^2}{u^2(1-1/r)^2+\pi^2}.
\]
The exact pole ratio satisfies
\begin{equation}\label{eq:negative-exact-ratio}
 \log R_q=-(n+1)D_q+\sum_i g_q(r_i).
\end{equation}
For $p=r/(1+r)$ and $w=p(1-p)$, the needed kernel bounds are
\begin{align}
 g_3(r)&\le D_3p+4\log3\,w,\label{eq:kernel-three}\\
 g_q(r)&\le D_qp+[4\log q+(q-1)/2]w,\label{eq:kernel-variance}\\
 g_q(r)&\le\max(q,4\log q)\,p.\label{eq:kernel-mean}
\end{align}
We give proofs to fix the constants.

For $r\ge1$, put $b=1/r\in[0,1]$ and $L=\log q$. Convexity of
$z\mapsto\tfrac12\log[(u^2z+q^2\pi^2)/(u^2z+\pi^2)]$ gives
\[
 g_q(r)\le D_q+(2b-b^2)(L-D_q)
 \le\frac{D_q}{1+b}+\frac{4Lb}{(1+b)^2}.
\]
The last comparison is affine in $D_q\in[0,L]$; at its endpoints it
reduces to $(1-b)^2(b+2)\ge0$ and $b(3-b)\ge0$.
This proves the right half of \eqref{eq:kernel-three} and
\eqref{eq:kernel-variance}.
For $r\le1$, $g_3$ decreases as $u$ increases, so it suffices to take
$u=\pi$. Put
\[
 N_r=1-2r+10r^2,\quad Q_r=1-2r+2r^2,\quad
 H(r)=\frac{(1+r)^2}{r}\log\frac{N_r}{Q_r}.
\]
Differentiation gives
\[
 H'(r)=(1-r)(1+r)
 \left[-\frac1{r^2}\log\frac{N_r}{Q_r}
       +\frac{16(1+r)}{N_rQ_r}\right]\ge0.
\]
For the inequality, use $\log z\le(z-1)/\sqrt z$ and
\[
 4(1+r)^2-N_rQ_r
 =3+12r(1-r)+4r^3(6-5r)\ge0.
\]
Since $H(1)=8\log3$, this proves $g_3\le4\log3\,w$ on the left.
For general $q$, put $y=(1-r)/r\ge0$ and $k=u/\pi\ge1$. Then
\[
 \partial_q(g_q/w)
 =\frac{q[y+3+1/(y+1)]}{k^2y^2+q^2}
 \le\frac{q(y+4)}{y^2+q^2}\le\frac12+\frac4q.
\]
The last step follows from $2qy\le y^2+q^2$. Integration from $1$
to $q$ proves \eqref{eq:kernel-variance}. Finally, if $r\ge1/3$,
then $p\ge1/4$ and $g_q\le\log q$. If $r\le1/3$, then
$y=1/r-1\ge2$, and $\log(1+t^2)\le t$ gives
\[
 \frac{g_q}{p}\le\frac{y+2}{2}\log(1+q^2/y^2)
 \le\frac{q(y+2)}{2y}\le q.
\]
This proves \eqref{eq:kernel-mean}.

Put $C_q=4\log q+(q-1)/2$ and $M_q=\max(q,4\log q)$.
Equations~\eqref{eq:bulk-mode-mean}, \eqref{eq:negative-exact-ratio}
and Lemma~\ref{lem:variance} yield
\begin{align}
 \log R_3&\le-n\left(1-\frac1{Y-\log2}\right)D_3
                    +\frac{32n\log3}{Y^2},\label{eq:negative-three-bound}\\
 \log R_q&\le-n\left(1-\frac1{Y-\log2}\right)D_q
                    +\frac{8nC_q}{Y^2},\label{eq:negative-var-bound}\\
 \log R_q&\le-(n+1)D_q+M_q\left(\frac n{Y-\log2}+1\right).
                    \label{eq:negative-mean-bound}
\end{align}

For $0\le u\le\pi$, use $B=(1+e^u)/\pi$ instead. Bound
\eqref{eq:bulk-mode-mean} remains valid: at the midpoint the
consecutive probability ratio is less than one because
$B\ge2/\pi$ and $n\ge250$; log-concavity therefore places the mode
before the midpoint. The kernel becomes
\[
 \frac12\log\frac{(u/\pi-1/r)^2+q^2}{(u/\pi-1/r)^2+1},
\]
and \eqref{eq:kernel-mean} still holds. Its proof for $r\le1/3$ uses
$|u/\pi-1/r|\ge1/r-1\ge2$; its other case is unchanged.
Thus \eqref{eq:negative-mean-bound} applies in this range as well.

\subsection{Summation of the negative-side pole bounds}

On $0\le u\le\pi$, the rational interval checks specified below
show that the bounds for $R_3,R_5,R_7,R_9$ are all less than $e^{-12}$.
For $Q=9$, the exponent on the right of
\eqref{eq:universal-pole-tail}, before its factor $Q/2$, is also less
than $-12$. Hence the whole sum is less than $.001$.

On $\pi\le u\le50$, the same checks give
$R_3<.79$, $R_5<e^{-4}<.02$, and bounds $e^{-12}$ for the remaining
finite poles and the tail exponent. There are at most $23$ finite
terms with $q\ge7$, and $Q/2\le25.5$, so the total is less than $.811$.

It remains to treat $50\le u\le4\sqrt n$. Here
\[
 Y\ge u+1.37,\quad Y-\log2\ge u+.67,\quad
 1-1/(Y-\log2)>.96,\quad n/u^2\ge1/16.
\]
The first inequality follows from
$n/u\ge\sqrt{250}/4>e^{1.37}$. Using
$\log(1+x)\ge2x/(2+x)$, the bound for $q=3$ reduces to
\begin{equation}\label{eq:negative-rational-polynomial}
 \left(1-\frac1{u+.67}\right)
       \frac{39.44u^2}{u^2+49.35}
 -\frac{35.156u^2}{(u+1.37)^2}\ge3.9.
\end{equation}
After multiplication by the positive common denominator, the
numerator has positive coefficients in $u-50$. Thus
$R_3\le e^{-3.9/16}<.79$.
For $q=5$,
\[
 -\log R_5\ge\frac n{u^2}
 \left\{\frac{.96\cdot9.86\cdot24}
 {2+9.87\cdot26/2500}-8C_5\right\}>2.53,
\]
so $R_5<.08$.
Let $Q$ be the smallest odd integer at least $u$. For $7\le q\le Q$,
$u\ge\max(50,q-2)$, and the following rational-polynomial comparison
holds:
\[
 \frac{.96\cdot9.86(q^2-1)}
 {2+9.87(q^2+1)/\max(50,q-2)^2}
 -32\log q-4(q-1)\ge16q.
\]
For completeness, replace $\log q$ by its tangent bound
$1.946+(q-7)/7$. With
\[
 A=.96\cdot9.86,\quad D=9.87,\quad
 L(q)=\frac{172}{7}q+26.272,
\]
the required numerators are
\[
 2500A(q^2-1)-L(q)[5000+D(q^2+1)]\quad(7\le q\le52)
\]
and
\[
 A(q^2-1)(q-2)^2
 -L(q)[2(q-2)^2+D(q^2+1)]\quad(q\ge52).
\]
The first has positive Bernstein coefficients on $[7,52]$ and the
second has positive coefficients in $q-52$. Therefore $R_q\le e^{-q}$,
and the sum from $q=7$ through $Q$ is less than $.0011$.

At $Q$, one has $D_Q>1.18$ and
\[
 \frac{8C_Q}{Y^2}
 \le\frac{32\log(u+2)+4(u+1)}{(u+1.37)^2}<.3.
\]
The last expression decreases for $u\ge50$, and its value at $50$
is less than $.3$. Hence $R_Q<e^{-.8n}$. The extra exponent in
\eqref{eq:universal-pole-tail} is at most $.051n$, and
$Q/2<3\sqrt n$. The infinite tail is less than $.001$.
Altogether,
\[
 \sum_{q\ge3,\ q\text{ odd}}R_q
 <.79+.08+.0011+.001<.88.
\]
This establishes all the quantitative bounds in Lemma~\ref{lem:bulk}.
At a phase node, the nearest conjugate pair contributes
$2C_s(-1)^j|\mathcal R_{n,s}(s+i\pi)|$; the remaining pairs have
strictly smaller total modulus. This proves its sign conclusion.

\subsection{Exact verification of the scalar inequalities}

We specify the finite checks so that the numerical constants in the
proof are reproducible independently of root computations. Every
terminating decimal in this section denotes the corresponding rational
number. The elementary bounds $3.14159<\pi<355/113$ suffice.
For $x\ge0$ and $x<m+2$, use
\[
 E_m(x)=\sum_{k=0}^m\frac{x^k}{k!},\qquad
 E_m(x)\le e^x\le
 E_m(x)+\frac{x^{m+1}}{(m+1)!}\frac1{1-x/(m+2)}.
\]
Take $m=100$ in the variance checks, $m=80$ in the bounded-interval
pole checks, and $m=220$ for the final tail exponent at $n=250$.
For logarithms, first reduce the argument to $[1,2]$ by powers of
$2$ and use, with $z=(x-1)/(x+1)$,
\[
 2\sum_{k=0}^{39}\frac{z^{2k+1}}{2k+1}
 \le\log x\le
 2\sum_{k=0}^{39}\frac{z^{2k+1}}{2k+1}
 +\frac{2z^{81}}{81(1-z^2)}.
\]
Reciprocation handles arguments below one.
For $I_0(v)$ on $[.7,1]$, retain its first $25$ terms and bound its
remaining terms by the geometric ratio $(v^2/4)/26^2$.

The following explicit upper bounds control the two series used in
the variance proof. Put $r=e^{-8t}$ and define
\begin{align*}
 \overline S(t)&=1+4e^{-2t}+9e^{-6t}
 +e^{-12t}\left\{\frac{16}{1-r}
 +\frac{8r}{(1-r)^2}+\frac{r(1+r)}{(1-r)^3}\right\},\\
 \overline T(t)&=e^{-2t}+4e^{-6t}
 +e^{-12t}\left\{\frac9{1-r}
 +\frac{6r}{(1-r)^2}+\frac{r(1+r)}{(1-r)^3}\right\}.
\end{align*}
They follow from $j(j-1)\ge12+8(j-4)$ for $j\ge4$.
Replacing exponentials by the rational bounds above makes every entry
of Table~\ref{tab:variance-intervals} a rational comparison. In its last
row, $\overline\Delta$ is obtained from $\Delta$ by replacing $S,T$
with $\overline S,\overline T$.

\begin{table}[ht]
\centering
\caption{Finite scalar checks for Lemma~\ref{lem:variance}.}
\label{tab:variance-intervals}
\begin{tabular}{p{.30\textwidth}p{.62\textwidth}}
\toprule
Intervals or parameters & Verified inequality\\
\midrule
$[a,b]=[.7+i/1000,.7+(i+1)/1000]$, $0\le i<300$
& $b[e^{-a}I_0(a)]^2<.47^2$\\[3pt]
$[a,b]$ one of $(2,3)$, $(3,4)$, $(4,6)$, $(6,8)$, $(8,10)$, $(10,11)$
& $\{.47[\overline S(a/7)+e^{-a/2}\overline S(a/3)+.003]\}^2<(8/b)^3$\\[3pt]
$[a,b]=[i/2,(i+1)/2]$, $22\le i<34$
& $G=8/b-1/4-\overline\Delta(a)>0$ and $\overline\Delta(a)<G^2$\\
\bottomrule
\end{tabular}
\end{table}

Here are the bounded-interval pole checks in equally explicit form.
For a box $[a,b]$, use
\[
 D_q^-:=\frac12\log\left(1+
       \frac{(q^2-1)9.86}{b^2+9.87}\right),\qquad
 \ell_2=.7>\log2.
\]
On $0\le u\le\pi$, take $[a,b]=[i/20,(i+1)/20]$,
$0\le i<63$, and
$Y^- =\log[250(1+e^a)/(355/113)]$.
For $q=3,5,7,9$, set
\[
 f_q=-D_q^-+\frac{M_q}{Y^--\ell_2},\qquad
 c_q=M_q-D_q^-.
\]
The checks give $f_q<0$ and $250f_q+c_q<-12$.
For $q=9$, they also give
\[
 f_q+\tau<0,\qquad 250(f_q+\tau)+c_q+\tau<-12,
 \quad \tau=\frac{b^2}{2q^2\,9.86}.
\]

On $\pi\le u\le50$, take
$a=3.14+i/10$, $b=\min(a+.1,50)$, $0\le i<469$, and
$Y^-=a+\log(250/b)$. Bound $q=3$ by
\eqref{eq:negative-three-bound}, using $D_3^-$, $Y^-$ and $.7$;
the resulting exponent at $n=250$ is less than $-39/160$.
Let $Q$ be the least odd integer at least $\max(5,b)$.
For each $q=5,7,\ldots,Q$, use the smaller of the mean bound and the
variance bound, with the same substitutions. The exponent is less
than $-4$ for $q=5$ and less than $-12$ for $q\ge7$.
For $q=Q$, adding $(n+1)\tau$ still gives an exponent below $-12$.
In each case the coefficient of $n$ is strictly negative and decreases
as $Y$ increases. Since $Y_n=Y_{250}+\log(n/250)$, the $n=250$ check
therefore proves the assertion for every $n\ge250$.

The remaining checks are polynomial positivity in
\eqref{eq:negative-rational-polynomial} and in the two displayed
polynomials in $q$, together with their scalar endpoint comparisons.
For the cubic on $[7,52]$, substitute $q=7+45t$ and convert its power
coefficients $a_k$ to Bernstein coefficients
$\sum_{k\le i}a_k\binom ik/\binom3k$; all four are positive.
For the unbounded intervals, direct expansion in $u-50$ or $q-52$
has strictly positive coefficients. Thus the finite verification uses
rational arithmetic throughout.
